% ============================================================
% CAPÍTULO 3
% ============================================================
\capitulo{Técnicas de Ciencia de Datos y Visión por Computadora}

Este capítulo reúne el marco teórico y el estado del arte del trabajo. Comienza con el relevamiento de los antecedentes en reconocimiento automático de dolor y emoción animal, y precisa el vacío que el trabajo busca cubrir. Luego presenta las técnicas que componen el sistema: estimación de pose, representación visual con redes convolucionales, rasgos de pose con bosques aleatorios, aprendizaje de múltiples instancias, combinación de clasificadores y, por último, los recaudos de validación y las métricas que gobiernan la comparación de modelos. La secuencia concreta de procesamiento y las decisiones de implementación se describen en el Capítulo~4.

\section{Reconocimiento Automático de Dolor y Emoción Animal con Visión por Computadora}
\textcite{broome2023} ofrecen la revisión de referencia sobre el reconocimiento por visión por computadora del dolor y de las emociones animales. Su título, «Going deeper than tracking», resume la dirección del campo: de seguir al animal a inferir su estado. La revisión señala dificultades que reaparecen en la literatura posterior: pocos datos etiquetados, etiquetas de definición discutible y generalización incierta a individuos no vistos.

En caballos, los primeros trabajos mostraron la viabilidad de reconocer expresiones de dolor mediante aprendizaje automático \parencite{andersen2021} y mediante redes convolucionales aplicadas a imágenes \parencite{lencioni2021}. \textcite{hummel2020} abordaron la detección automática de dolor en caras de caballos y burros, y \textcite{pessanha2023} compararon enfoques basados en imágenes faciales para la evaluación del dolor equino. Con video, \textcite{broome2019} encontraron que la dinámica temporal es importante para reconocer el dolor equino, mientras que trabajos posteriores emplearon representaciones de pose aprendidas sin supervisión \parencite{rashid2022} y estudiaron la transferencia entre dominios de dolor \parencite{broome2022}. Fuera del equino, el reconocimiento automático de dolor en gatos mostró también ser factible \parencite{feighelstein2022}.

Los antecedentes sobre estados afectivos distintos del dolor son menos numerosos. \textcite{corujo2021} aplicaron redes convolucionales al reconocimiento de emociones en caballos a partir de imágenes, y \textcite{feighelstein2024} abordaron el reconocimiento de estados emocionales de caballos a partir de la expresión facial. Una línea que atraviesa varios de estos trabajos es la comparación entre entradas basadas en puntos de referencia geométricos y entradas basadas en la apariencia de la imagen; esa dicotomía es la que este trabajo aprovecha al combinar rasgos de pose con rasgos de una red preentrenada.

\subsection{Vacío que Cubre este Trabajo}
Del relevamiento surgen tres rasgos del estado del arte. Primero, el grueso del trabajo automático en caballos se ocupa del dolor y no de categorías afectivas amplias, y los antecedentes sobre emoción son pocos \parencite{corujo2021,feighelstein2024}. Segundo, los trabajos con video suelen emplear secuencias más largas y modelos temporales \parencite{broome2019}. Tercero, la mayoría de los estudios dispone de pocos individuos, de modo que la forma de evaluar determina cuánto puede creerse el resultado \parencite{varoquaux2018,roberts2017}.

Este trabajo cubre el siguiente hueco: la clasificación de tres clases afectivas (\clase{Sereno}, \clase{Interesado} y \clase{Disgustado}) en videos cortos, sobre un conjunto de datos propio, combinando la pose facial (modelo SuperAnimal y rasgos de estilo SimBA) con la apariencia (red residual preentrenada), y con una evaluación explícitamente rigurosa: validación agrupada por caballo, reserva de caballos de un solo uso, intervalos por \textit{bootstrap} por caballo y validación dejando afuera una sesión de filmación. No se pretende superar a los antecedentes en dolor, que resuelven un problema distinto con otros datos, sino estudiar qué parte de la señal afectiva puede recuperarse en condiciones no controladas y cuán confiable es esa estimación.

\section{Estimación de Pose y Análisis Automático de la Conducta}
DeepLabCut popularizó la estimación de pose sin marcadores mediante redes profundas entrenadas con pocos fotogramas anotados \parencite{mathis2018}. Su aplicación a caballos, incluida la robustez fuera de dominio con preentrenamiento, fue estudiada por \textcite{mathis2021}, y las revisiones del área subrayan que los puntos clave son una entrada intermedia para análisis posteriores \parencite{mathis2020}. Los modelos SuperAnimal reducen la necesidad de anotar manualmente en numerosas especies \parencite{ye2024}; este trabajo usa el modelo preentrenado SuperAnimal-Quadruped, que devuelve 39 puntos con una confianza para cada uno, sin anotar datos propios. Un fotograma se considera válido cuando el ojo y el hocico superan una confianza de 0,50.

A partir de los puntos clave se desarrollaron herramientas para clasificar conductas. SimBA construye rasgos geométricos a partir de la pose y los entrega a clasificadores supervisados, con herramientas de interpretación \parencite{goodwin2024}; es la inspiración directa de los \PoseRasgos{} rasgos «de estilo SimBA» de este trabajo. Existen también alternativas no supervisadas, como B-SOiD \parencite{hsu2021} y Keypoint-MoSeq \parencite{weinreb2024}, que descubren motivos de conducta sin etiquetas. No se emplean aquí porque el problema tiene categorías afectivas definidas por una experta. Conviene señalar que estas herramientas se desarrollaron principalmente para animales de laboratorio con cámaras fijas, de modo que su uso sobre videos de caballos en condiciones no controladas constituye un aporte y no una aplicación rutinaria.

\section{Normalización Espacial de la Cara}
Con el ojo y el hocico de cada fotograma se estima una transformación de similitud que corrige rotación, escala y traslación y lleva la cara a una posición canónica. El recorte resultante mide $224 \times 343$ píxeles; la altura adicional respecto del formato cuadrado habitual conserva las orejas. Los bordes se rellenan replicando los píxeles vecinos para no introducir bandas negras cuyo tamaño dependa del encuadre. El objetivo es reducir la posibilidad de que el modelo aprenda diferencias de encuadre en lugar de diferencias de expresión.

\section{Representación Visual: Redes Convolucionales y Aprendizaje por Transferencia}
Las redes convolucionales extraen representaciones jerárquicas de las imágenes, y las redes residuales permiten entrenar arquitecturas profundas mediante conexiones de atajo entre capas \parencite{he2016}. Con una cantidad reducida de datos, como la de este trabajo, resulta preferible reutilizar una red entrenada sobre una colección grande, como ImageNet \parencite{deng2009}, porque las capas convolucionales aprenden representaciones generales que se transfieren a tareas nuevas \parencite{yosinski2014}.

Se utiliza una ResNet-18 preentrenada como extractor de características: se elimina su clasificador y se conserva la salida del \textit{pooling} global, un vector de 512 dimensiones por fotograma. La red se mantiene congelada, es decir, sus parámetros no se actualizan. Con \CaballosTotal{} caballos, entrenar sus 11 millones de parámetros favorecería la memorización de los individuos; congelada, permite además precalcular las características una sola vez. A partir de esta representación se pueden construir clasificadores de diversa complejidad.

\subsection{Línea de Base Estática}
Una línea de base es el punto de comparación que permite saber si la complejidad adicional de un modelo se justifica. Aquí se usa una regresión logística multinomial sobre el promedio de las representaciones visuales de los fotogramas de un clip. Ese modelo descarta por construcción el orden temporal y la ponderación de instancias, de modo que cualquier ventaja de un modelo más elaborado sobre él mide cuánto aportan esos elementos. La regresión logística admite regularización, lo que la hace apropiada para vectores de 512 dimensiones con pocos cientos de ejemplos \parencite{pedregosa2011}.

\section{Rasgos de Pose y Bosques Aleatorios}
Para cada ventana de \Fotogramas{} fotogramas se calculan rasgos geométricos a partir de los puntos clave faciales, expresados respecto de una referencia ojo--hocico y escalados por la distancia entre ambos, de modo que sean aproximadamente invariantes al tamaño y a la posición de la cara. Siguiendo el enfoque de SimBA \parencite{goodwin2024}, se resumen estadísticamente a lo largo de la ventana, lo que da un vector tabular de \PoseRasgos{} rasgos por clip.

Para datos tabulares con pocas observaciones y rasgos correlacionados, el bosque aleatorio (\textit{random forest}) es un clasificador robusto \parencite{breiman2001}: promedia muchos árboles de decisión entrenados sobre muestras con reemplazo y subconjuntos aleatorios de rasgos, lo que reduce la varianza respecto de un único árbol. Su interpretación puede complementarse con valores SHAP, que asignan a cada rasgo una contribución a la predicción con propiedades axiomáticas \parencite{lundberg2017}. En este trabajo esa auditoría de los rasgos permitió contrastar si el modelo se apoyaba en la expresión o en características de la filmación.

\section{Modelos Sensibles e Invariantes al Orden Temporal}
Un clip de video aporta varios fotogramas, pero eso no implica que un modelo aproveche su orden. Conviene distinguir dos familias. Los modelos invariantes al orden tratan los fotogramas como un conjunto: votan entre clasificaciones individuales, promedian las representaciones, o las ponderan con atención sin información de posición. Los modelos sensibles al orden, en cambio, incorporan la secuencia: redes recurrentes como la GRU \parencite{cho2014} y la LSTM \parencite{hochreiter1997}, o mecanismos de atención con codificación posicional. La importancia de la dinámica temporal se ha documentado para el dolor equino con secuencias largas \parencite{broome2019}, y plantea una hipótesis natural para el problema afectivo: que modelar el orden de los fotogramas mejore la clasificación.

Contrastar esa hipótesis exige comparar un modelo sensible al orden con una alternativa equivalente que no lo use, de modo que cualquier diferencia pueda atribuirse al orden y no a otra diferencia de arquitectura. Esa lógica de comparación se retoma en el Capítulo~4; aquí interesa fijar que un modelo que procesa varios fotogramas no es, por esa sola razón, un modelo temporal, y que por eso este trabajo evita calificarlo como espacio-temporal salvo que la evidencia lo respalde.

\section{Aprendizaje de Múltiples Instancias con Atención}
En el aprendizaje de múltiples instancias (MIL, por sus siglas en inglés), cada ejemplo es una bolsa de instancias con una única etiqueta \parencite{dietterich1997}. La variante con atención aprende un peso por instancia y forma la representación de la bolsa como un promedio ponderado \parencite{ilse2018}. Dada una bolsa de instancias con representaciones $\mathbf{h}_1,\dots,\mathbf{h}_K$, la representación de la bolsa es
\begin{equation}
\mathbf{z} = \sum_{k=1}^{K} a_k \mathbf{h}_k, \qquad
a_k = \frac{\exp\{\mathbf{w}^{\top}\tanh(\mathbf{V}\mathbf{h}_k)\}}{\sum_{j=1}^{K}\exp\{\mathbf{w}^{\top}\tanh(\mathbf{V}\mathbf{h}_j)\}},
\label{eq:c3-atencion}
\end{equation}
donde $\mathbf{V}$ y $\mathbf{w}$ son parámetros aprendidos. Aquí cada clip es una bolsa de \Fotogramas{} fotogramas. Como los pesos de la ecuación~(\ref{eq:c3-atencion}) no dependen de la posición, la atención sin codificación posicional es, por construcción, invariante al orden de las instancias. Este modelo fue la propuesta original del trabajo y se evalúa en el Capítulo~5 como uno de los cuatro candidatos. Ningún candidato de la comparación final incorpora el orden de los fotogramas; las redes recurrentes GRU \parencite{cho2014} y LSTM \parencite{hochreiter1997} se mencionan como referencia del estado del arte y como parte de exploraciones previas descartadas (Anexo~C), no como candidatos de la comparación.

\section{Combinación de Clasificadores y Fusión Tardía}
Cuando se dispone de dos representaciones distintas del mismo ejemplo, una alternativa a concatenarlas es entrenar un clasificador por representación y combinar sus salidas. \textcite{kittler1998} estudian un marco común para esa combinación y analizan reglas como el producto, la suma, el máximo, el mínimo, la mediana y el voto. Este trabajo emplea la regla de la suma con pesos, en una fusión tardía de un clasificador sobre la pose y otro sobre la apariencia visual. Si $\mathbf{p}_{\mathrm{pose}}(x)$ y $\mathbf{p}_{\mathrm{vis}}(x)$ son los vectores de probabilidades por clase de cada clasificador para el clip $x$, la probabilidad fusionada es
\begin{equation}
\mathbf{p}(x) = w\,\mathbf{p}_{\mathrm{pose}}(x) + (1-w)\,\mathbf{p}_{\mathrm{vis}}(x), \qquad
\hat{y}(x) = \arg\max_{c}\; p_c(x),
\label{eq:c3-fusion}
\end{equation}
con $w \in [0,1]$ el peso de la pose. El peso se elige por validación interna, usando únicamente datos de entrenamiento, y no sobre las particiones de prueba; en el Capítulo~4 se detalla el procedimiento y el valor resultante, $w = \PesoPose{}$.

\section{Validación en Conjuntos Pequeños con Estructura de Grupos}
La forma de evaluar condiciona la credibilidad del resultado tanto como el modelo. Este trabajo adopta varios recaudos que la literatura señala para conjuntos pequeños y agrupados.

\subsection{Agrupamiento por Individuo y por Sesión}
Si clips del mismo caballo aparecen en entrenamiento y en prueba, el modelo puede reconocer al animal en lugar de la expresión, una forma de fuga de información que infla las métricas \parencite{kaufman2012}. Cuando los datos tienen estructura jerárquica o de agrupamiento, la validación cruzada debe respetarla \parencite{roberts2017}. Por eso todas las particiones agrupan los clips por caballo: ningún animal aparece a la vez en entrenamiento y en prueba. Como los caballos de una misma sesión comparten lugar, fecha y condiciones de filmación, se agrega una validación que deja afuera una sesión completa por vez. La fuga de información no es una curiosidad metodológica: se ha documentado como una causa extendida de problemas de reproducibilidad en la ciencia basada en aprendizaje automático \parencite{kapoor2023}.

\subsection{Muestras Pequeñas y Sobreajuste en la Selección de Modelos}
Con pocas observaciones, las estimaciones de validación cruzada tienen un error considerable \parencite{varoquaux2018}. Por eso se repite la validación \Repeticiones{} veces con \Folds{} particiones por caballo y se informa la variabilidad entre repeticiones. Además, elegir entre modelos sobre los mismos datos con los que se estima su desempeño introduce un sesgo de selección: el mejor de varios candidatos tiende a parecer mejor de lo que es \parencite{cawley2010}. Para contenerlo, la comparación entre candidatos se fija antes de medir, y un grupo de caballos (\CaballosReserva{} de \CaballosTotal{}) se reserva y se usa una sola vez, sobre el candidato ya elegido.

\subsection{Atajos y Confusores}
Los modelos de aprendizaje profundo pueden apoyarse en atajos, es decir, en regularidades del conjunto que correlacionan con la etiqueta sin reflejar la señal de interés \parencite{geirhos2020}. En este problema la asociación entre clase y contexto de filmación (Capítulo~2) es un candidato natural. La forma de detectarlo es evaluar el desempeño también sobre un subconjunto en que el atajo no actúa, y contrastar el resultado con el de la evaluación completa.

\subsection{Intervalos por Bootstrap Agrupado}
La incertidumbre de las métricas y de las diferencias entre modelos se estima con \textit{bootstrap} \parencite{efron1993}. Como los clips de un mismo caballo no son independientes, se remuestrean caballos y no clips; las diferencias entre modelos se calculan de manera pareada sobre los mismos remuestreos, con \ReplicasBootstrap{} réplicas.

\section{Métricas para Clases Desbalanceadas}
Las clases del conjunto no tienen el mismo tamaño, por lo que la exactitud puede ocultar un desempeño pobre en las clases menos frecuentes. La métrica principal es el \fmacro{}, el promedio no ponderado del \textit{F}\textsubscript{1} de cada clase, que otorga el mismo peso a cada una \parencite{sokolova2009}. Para la clase $c$, con precisión $P_c$ y sensibilidad $R_c$,
\begin{equation}
F_{1,c} = \frac{2\,P_c\,R_c}{P_c + R_c}, \qquad
F_{1}^{\mathrm{macro}} = \frac{1}{C}\sum_{c=1}^{C} F_{1,c},
\label{eq:c3-f1}
\end{equation}
donde $C = 3$ es el número de clases. Existen variantes del promedio macro (por ejemplo, calcular la \textit{F}\textsubscript{1} a partir de la precisión y la sensibilidad promedio), que pueden dar valores distintos \parencite{opitz2019}; aquí se usa siempre la definición de la ecuación~(\ref{eq:c3-f1}). Se informan también la exactitud, la precisión, la sensibilidad y la matriz de confusión, y los resultados por clase, para no esconder diferencias entre ellas.
